% Báo cáo thiết kế - Trò chơi "Phi đội không gian"
% Biên dịch: pdflatex report.tex (chạy 2 lần để có mục lục)
% Ảnh sơ đồ: thư mục images/ cạnh file .tex (Overleaf) hoặc ../design/diagrams/images/ (repo).
\documentclass[12pt,a4paper]{article}

\usepackage[utf8]{inputenc}
\usepackage[T5]{fontenc}
\usepackage[vietnamese]{babel}
\usepackage[margin=2.2cm]{geometry}
\usepackage{graphicx}
\usepackage{booktabs}
\usepackage{tabularx}
\usepackage{longtable}
\usepackage{array}
\usepackage{enumitem}
\usepackage{xcolor}
\usepackage{listings}
\usepackage{float}
\usepackage{caption}
\usepackage{seqsplit}
\usepackage{collcell}
\usepackage[hidelinks]{hyperref}

\graphicspath{{images/}{../design/diagrams/images/}}
\setlist{itemsep=1pt,topsep=3pt}
\setlength{\parskip}{4pt}
\renewcommand{\arraystretch}{1.15}

\lstset{basicstyle=\ttfamily\small, columns=fullflexible, frame=single,
        breaklines=true, keepspaces=true, backgroundcolor=\color{gray!6}}

\newcommand{\cmd}[1]{\texttt{\seqsplit{#1}}}
\newcolumntype{L}{>{\raggedright\arraybackslash}X}
\newcolumntype{P}[1]{>{\raggedright\arraybackslash}p{#1}}
% Ô ví dụ message: chữ máy nhỏ, xuống dòng thủ công bằng \newline.
\newcolumntype{E}[1]{>{\raggedright\arraybackslash\ttfamily\footnotesize}p{#1}}
% Ô JSON: cho phép ngắt ở mọi ký tự để request/response không tràn sang cột bên cạnh.
\newcommand{\jsoncell}[1]{{\raggedright\arraybackslash\ttfamily\footnotesize\seqsplit{#1}}}
\newcolumntype{J}[1]{>{\collectcell\jsoncell}p{#1}<{\endcollectcell}}
% Gạch đầu dòng gọn trong ô bảng.
\newcommand{\bul}{\textbullet~}

% Chèn sơ đồ; nếu chưa render ảnh thì hiện khung báo thiếu thay vì lỗi biên dịch.
\newcommand{\diagram}[3]{%
  \begin{figure}[H]\centering
  \IfFileExists{images/#1.png}%
    {\includegraphics[width=#3\linewidth]{images/#1.png}}%
    {\IfFileExists{../design/diagrams/images/#1.png}%
      {\includegraphics[width=#3\linewidth]{../design/diagrams/images/#1.png}}%
      {\fbox{\parbox{0.8\linewidth}{\centering\small Chưa có ảnh \texttt{#1.png}.\\
         Chạy \texttt{python3 -I Project/tools/render\_diagrams.py} để render.}}}}%
  \caption{#2}\label{fig:#1}
  \end{figure}}

\title{\textbf{BÁO CÁO THIẾT KẾ}\\[4pt]
       \large Ứng dụng trò chơi ``Phi đội không gian''\\[2pt]
       \normalsize Môn Lập trình mạng -- IT4062}
\author{Đặng Hoàng Sơn -- 20235994\\ Ngô Đức Huy -- 20235946}


\begin{document}
\maketitle
\tableofcontents
\newpage

% =====================================================================
\section{Giới thiệu}

\textbf{Đề tài.} Xây dựng trò chơi nhiều người chơi theo mô hình client--server trên TCP. Người chơi lập đội tối đa 3 người, thách đấu đội khác và điều khiển phi thuyền tấn công đối phương. Đội nào bị hạ hết phi thuyền trước thì thua. Coin để mua trang bị kiếm được bằng cách trả lời câu hỏi trắc nghiệm trong các rương kho báu.

\textbf{Mục tiêu thiết kế.}
\begin{itemize}
  \item Đáp ứng đầy đủ 23 chức năng của barem (39 điểm).
  \item Server chịu được nhiều kết nối và nhiều trận chạy cùng lúc.
  \item Server là bên quyết định duy nhất: client không tự tính luật chơi.
  \item Biên dịch và chạy được trên Ubuntu 20.04, GCC 9.4, chuẩn \cmd{-std=c11}.
\end{itemize}

% =====================================================================
\section{Đặc tả trò chơi}

\subsection{Đội và trận đấu}
\begin{itemize}
  \item Đội hoạt động như \emph{party} khi ghép trận: chỉ tồn tại trong RAM của server, gồm những người đang online. Mỗi đội có 3 slot cố định (0, 1, 2), tối đa 3 người.
  \item Thành viên mất kết nối khi đội không ở trong trận thì tự rời đội. Mất kết nối trong trận thì vẫn ở trong đội tới hết trận để còn vào lại được. Chỉ mời được người đang online.
  \item Hết trận, đội quay về room với các thành viên còn online, giữ nguyên đội trưởng, và thách đấu tiếp được ngay.
  \item Chỉ 2 đội cùng số người mới thách đấu được nhau. Đội đang ở sảnh chờ hoặc trong trận thì bị khóa mọi thao tác đội, kể cả với admin.
  \item Trong sảnh chờ, đội trưởng đội chủ (người tạo phòng) chỉnh thời gian lượt (mặc định 30 giây), thời gian rương (30 giây) và coin ban đầu (0). Đội khách sẵn sàng thì đội chủ mới bắt đầu được.
  \item Mỗi trận reset toàn bộ: coin, trang bị, HP, đạn. Coin tính riêng từng người.
\end{itemize}

\subsection{Tàu, vũ khí và giáp}
Tàu khởi đầu có HP 1000, Amor 0 và 1 pháo tự động với 50 viên đạn.

\begin{table}[H]\centering\small
\begin{tabularx}{\linewidth}{@{}l l l L@{}}
\toprule
\textbf{Trang bị} & \textbf{Giá (coin)} & \textbf{Thông số} & \textbf{Ghi chú} \\
\midrule
Pháo tự động & có sẵn & Damage 10/viên & Đúng 1 khẩu, không bỏ, không mua thêm. Thùng đạn 100 coin/50 viên \\
Pháo laze & 1000/khẩu & Damage 100/phát & Tối đa 3 khẩu (pháo tự động chiếm 1 trong 4 slot). Bộ nguồn 100 coin/10 phát \\
Tên lửa & 2000/quả & Damage 800 & Dùng 1 lần, không chiếm slot pháo, mang tối đa 4 quả \\
Giáp cơ bản & 1000 & Amor 500 & Tối đa 2 lớp, cộng dồn \\
Giáp tăng cường & 2000 & Amor 1500 & Lớp mua sau bị trừ trước \\
Sửa tàu & 1/HP & -- & Tối đa 1000 HP; tàu đã chết không sửa được \\
\bottomrule
\end{tabularx}
\caption{Bảng trang bị}
\end{table}

\begin{itemize}
  \item \textbf{Đạn:} mỗi khẩu có băng riêng; kho thùng đạn và bộ nguồn dùng chung cho cùng loại. Reload làm được bất cứ lúc nào, mỗi lần lấy 1 thùng hoặc 1 bộ, số đạn còn trong băng cũ bị bỏ. Băng hết thì tự reload nếu kho còn.
  \item \textbf{Giáp:} đang đủ 2 lớp mà mua thêm thì lớp còn ít Amor hơn bị bỏ. Lớp về 0 Amor thì bị gỡ.
  \item \textbf{Sát thương:} mỗi điểm Damage trừ 1 Amor, bắt đầu từ lớp ngoài; phần dư trừ vào HP. HP về 0 thì tàu bị hạ, không hồi sinh.
  \item \textbf{Tấn công:} chọn 1 mục tiêu, sau đó chọn phần pháo (chỉ 1 loại: tự động hoặc laze, chọn số khẩu), phần tên lửa (chọn số quả), hoặc cả hai. Không được trộn 2 loại pháo.
\end{itemize}

\subsection{Lượt và rương kho báu}
\begin{itemize}
  \item Thứ tự lượt A0 $\rightarrow$ B0 $\rightarrow$ A1 $\rightarrow$ B1 $\rightarrow$ A2 $\rightarrow$ B2, bỏ qua tàu đã chết. Trong lượt của mình: mua, sửa, reload, mở rương; bấm tấn công là hết lượt. Ngoài lượt chỉ được reload.
  \item Hết giờ lượt hoặc mất kết nối thì bị bỏ qua lượt.
  \item Mỗi lượt server thả 1 rương: đồng 60\% (100 coin), bạc 30\% (500), vàng 10\% (2000). Mở rương là không bắt buộc.
  \item Mở rương thì đồng hồ lượt tạm dừng và rương có đồng hồ riêng. Trong lúc đó chỉ được trả lời. Câu hỏi trắc nghiệm, trả lời 1 lần: đúng thì nhận coin, sai thì mất rương. Đồng đội chỉ được xem câu hỏi; đối phương không thấy.
  \item Đội bị hạ hết tàu thì thua. Cả đội offline liên tục 120 giây thì thua do bỏ cuộc.
\end{itemize}

% =====================================================================
\section{Kiến trúc hệ thống}

\subsection{Mô hình server: 1 luồng chính epoll + thread pool}
\begin{itemize}
  \item \textbf{Luồng chính} chạy \cmd{epoll\_wait}, nhận/gửi trên mọi socket non-blocking, và là luồng \emph{duy nhất} được đọc/sửa state (phiên, đội, sảnh chờ, trận, đồng hồ).
  \item \textbf{Thread pool} chỉ làm việc chậm: hash mật khẩu, đọc SQLite, đọc file log. \textbf{1 luồng DB writer} ghi DB và ghi log theo đúng thứ tự FIFO. Kết quả được trả về luồng chính qua hàng đợi và \cmd{eventfd}.
  \item \textbf{Đồng hồ} lượt, rương, idle, bỏ cuộc nằm trong một min-heap; mốc gần nhất làm timeout của \cmd{epoll\_wait}. Không dùng \cmd{sleep} hay luồng timer riêng.
\end{itemize}

\subsection{Các quyết định thiết kế khác}
\begin{table}[H]\centering\small
\begin{tabularx}{\linewidth}{@{}l L@{}}
\toprule
\textbf{Vấn đề} & \textbf{Lựa chọn} \\
\midrule
Định dạng message & JSON (thư viện cJSON) \\
Truyền dòng & Mỗi message 1 dòng JSON, kết thúc bằng \cmd{\textbackslash r\textbackslash n} \\
Đồng bộ client & Event kèm snapshot toàn bộ trạng thái trận \\
Client & Socket gắn vào main loop của GTK, 1 luồng \\
Lưu trữ & SQLite, chế độ WAL \\
\bottomrule
\end{tabularx}
\caption{Các quyết định kiến trúc}
\end{table}

\subsection{Lưu trữ}
\begin{itemize}
  \item SQLite chỉ có 2 bảng: \cmd{users} (tên đăng nhập, hash PBKDF2 kèm salt, role, trạng thái khóa) và \cmd{questions} (tier, câu hỏi, 4 lựa chọn, đáp án). Đội (party), phiên, sảnh chờ và trận đang chơi chỉ nằm trong RAM; kết quả trận được ghi vào log.
  \item Lúc khởi động, server nạp tài khoản (trừ hash mật khẩu) và câu hỏi lên RAM; luồng chính chỉ đọc từ RAM. Ghi theo kiểu \emph{RAM trước, DB sau} qua 1 luồng DB writer chạy FIFO, nên luồng chính không bao giờ chờ ổ đĩa. Riêng \cmd{REGISTER} phải chờ DB cấp \cmd{user\_id} nội bộ.
  \item Admin đầu tiên được tạo trên máy server bằng \cmd{./server --create-admin <username>}; câu hỏi soạn sẵn nạp từ \cmd{data/questions.txt}.
\end{itemize}

\subsection{Ghi log}
Server ghi lại các sự kiện quan trọng để quản trị viên theo dõi hoạt động và hỗ trợ tìm lỗi. Mỗi bản ghi nằm trên một dòng, gồm thời gian ISO-8601, mức độ, nhóm sự kiện, tên sự kiện và các cặp \cmd{khóa=giá\_trị} liên quan.

\begin{itemize}
  \item Các nhóm log chính gồm \cmd{NET}, \cmd{AUTH}, \cmd{TEAM}, \cmd{LOBBY}, \cmd{MATCH}, \cmd{CHEST}, \cmd{ADMIN}, \cmd{DB} và \cmd{SYS}.
  \item Các sự kiện cần ghi gồm kết nối/ngắt kết nối, đăng nhập, thao tác đội, bắt đầu/kết thúc trận, tấn công, mở rương, thao tác quản trị và lỗi cơ sở dữ liệu hoặc hệ thống.
  \item Luồng chính chỉ tạo bản ghi và đưa vào hàng đợi; luồng DB writer ghi ra file theo thứ tự FIFO. Nhờ đó thao tác ổ đĩa không chặn việc xử lý kết nối và luật chơi.
  \item Log không chứa mật khẩu, hash mật khẩu, salt hoặc nội dung nhạy cảm khác. Quản trị viên đọc log thông qua lệnh \cmd{ADMIN\_LOG}.
\end{itemize}

\begin{lstlisting}
2026-11-17T10:26:10.777+07:00 INFO MATCH ATTACK match=3 from=son to=minh damage=1000
\end{lstlisting}

% =====================================================================
\section{Giao thức}

\subsection{Đóng gói và cấu trúc message}
Mỗi message là 1 object JSON UTF-8 trên 1 dòng, kết thúc bằng \cmd{\textbackslash r\textbackslash n}, tối đa 8192 byte. Có 3 loại message:

\begin{table}[H]\centering\small
\begin{tabularx}{\linewidth}{@{}P{0.15\linewidth} P{0.2\linewidth} L E{0.33\linewidth}@{}}
\toprule
\textbf{Loại} & \textbf{Trường} & \textbf{Ý nghĩa} & \normalfont\small\textbf{Ví dụ message} \\
\midrule
Request \newline (C$\rightarrow$S) & \cmd{cmd}, tham số & \cmd{cmd} xác định lệnh cần thực hiện; mỗi kết nối chỉ có một request đang chờ response &
\{"cmd":"ATTACK",\newline "target":"minh",\newline "gun":\{"type":"LASER",\newline \ \ "count":2\},\newline "missiles":1\} \\
\midrule
Response \newline (S$\rightarrow$C) & \cmd{type}, \cmd{code}, \cmd{msg}, \cmd{data} & Đúng 1 response cho request đang chờ; \cmd{code} thuộc nhóm \cmd{1xxx} thành công, \cmd{2xxx} định dạng, \cmd{3xxx} tài khoản/phiên, \cmd{4xxx} phân quyền, \cmd{5xxx} đội, \cmd{6xxx} thách đấu/sảnh chờ, \cmd{7xxx} trong trận, \cmd{9xxx} lỗi server &
\{"type":"res",\newline "code":"1xxx","msg":"OK"\} \\
\midrule
Event \newline (S$\rightarrow$C) & \cmd{type}, \cmd{evt}, \cmd{data} & Server tự gửi; event trong trận kèm \cmd{state} &
\{"type":"evt",\newline "evt":"ATTACK",\newline "data":\{"attacker":"son",\newline \ \ "target":"minh",\newline \ \ "damage":1000,...\}\} \\
\bottomrule
\end{tabularx}
\caption{Ba loại message}
\end{table}

\subsection{Mã kết quả}
\begin{table}[H]\centering\small
\begin{tabularx}{\linewidth}{@{}l l L@{}}
\toprule
\textbf{Nhóm} & \textbf{Khoảng} & \textbf{Ví dụ} \\
\midrule
Thành công & 1xxx & OK \\
Định dạng & 2xxx & JSON hỏng, lệnh không tồn tại, thiếu tham số, message quá dài \\
Tài khoản, phiên & 3xxx & Chưa đăng nhập, sai mật khẩu, tài khoản bị khóa, trùng tên \\
Phân quyền & 4xxx & Không đủ quyền, không phải đội trưởng, không phải admin \\
Đội & 5xxx & Đội đầy, đã có đội, đội đang trong trận, người được mời đang offline \\
Thách đấu, sảnh chờ & 6xxx & Lệch số người, đội khách chưa sẵn sàng, setting sai \\
Trong trận & 7xxx & Chưa tới lượt, thiếu coin, hết slot, trộn pháo, đang mở rương \\
Server & 9xxx & Lỗi nội bộ hoặc chế độ chỉ đọc, server quá tải \\
\bottomrule
\end{tabularx}
\caption{Các nhóm mã kết quả}
\end{table}

\subsection{Chi tiết lệnh}

Mọi request có trường bắt buộc \cmd{cmd} (chuỗi). Tham số của lệnh nằm \textbf{cùng cấp} với \cmd{cmd}, không bọc trong \cmd{data}. Mọi response có \cmd{type="res"}, \cmd{code}, \cmd{msg}; trường \cmd{data} chỉ xuất hiện khi lệnh thành công và có dữ liệu cần trả về. Ký hiệu ``--'' trong bảng nghĩa là không có tham số hoặc không có \cmd{data}.

Các kiểu dùng chung:
\begin{itemize}
  \item \cmd{username}: chuỗi 3--20 ký tự, chỉ gồm chữ, số và dấu gạch dưới; \cmd{password}: chuỗi 6--64 ký tự.
  \item \cmd{team\_id}, \cmd{lobby\_id}, \cmd{question\_id}: số nguyên dương; \cmd{accept}, \cmd{ready}, \cmd{locked}: boolean JSON \cmd{true}/\cmd{false}.
  \item \cmd{role}: \cmd{"PLAYER"} hoặc \cmd{"ADMIN"}; trạng thái người chơi/đội: \cmd{"IDLE"}, \cmd{"LOBBY"} hoặc \cmd{"IN\_MATCH"}.
  \item \cmd{member}: object gồm \cmd{slot:int}, \cmd{username:string}, \cmd{is\_captain:bool}, \cmd{online:bool}.
  \item \cmd{settings}: object gồm \cmd{turn\_time:int} (10--120 giây), \cmd{chest\_time:int} (10--120 giây), \cmd{start\_coin:int} (0--100000).
  \item Các lệnh trả danh sách dài dùng \cmd{offset:int} (mặc định 0) và \cmd{limit:int} (mặc định 20, tối đa 50). Response trả thêm \cmd{has\_more:bool} và \cmd{next\_offset:int|null}.
\end{itemize}

Để mọi message luôn nhỏ hơn giới hạn 8192 byte, server có thể trả ít hơn \cmd{limit} phần tử nếu thêm phần tử tiếp theo làm response vượt giới hạn. Khi đó \cmd{has\_more=true} và client dùng \cmd{next\_offset} cho request kế tiếp. Quy tắc này áp dụng cho \cmd{TEAM\_LIST}, \cmd{ADMIN\_USER\_LIST}, \cmd{ADMIN\_LOG} và \cmd{ADMIN\_QUESTION\_LIST}. Thứ tự phân trang phải cố định: đội theo \cmd{team\_id} tăng dần, người dùng theo \cmd{username}, câu hỏi theo \cmd{question\_id}, log theo thứ tự dòng từ cũ đến mới. Nếu dữ liệu thay đổi trong lúc đang phân trang, client có thể tải lại từ \cmd{offset=0}.

\begin{longtable}{@{}P{0.18\linewidth} J{0.37\linewidth} J{0.37\linewidth}@{}}
\toprule
\textbf{Chức năng} & \multicolumn{1}{P{0.37\linewidth}}{\textbf{Request message}} & \multicolumn{1}{P{0.37\linewidth}}{\textbf{Response message}} \\
\midrule\endfirsthead
\toprule
\textbf{Chức năng} & \multicolumn{1}{P{0.37\linewidth}}{\textbf{Request message}} & \multicolumn{1}{P{0.37\linewidth}}{\textbf{Response message}} \\
\midrule\endhead

\normalfont Ping \newline \cmd{PING} &
\{"cmd":"PING"\} &
\{"type":"res","code":1000,"msg":"OK","data":\{"time":1763430116000\}\} \\
\midrule

\normalfont Đăng ký \newline \cmd{REGISTER} &
\{"cmd":"REGISTER",\newline "username":"son","password":"Abc123"\} &
\{"type":"res","code":1000,"msg":"OK"\} \\
\midrule

\normalfont Đăng nhập \newline \cmd{LOGIN} &
\{"cmd":"LOGIN",\newline "username":"son","password":"Abc123"\} &
\{"type":"res","code":1000,"msg":"OK",\newline "data":\{"username":"son","role":"PLAYER","team\_id":null\}\} \\
\midrule

\normalfont Đăng xuất \newline \cmd{LOGOUT} &
\{"cmd":"LOGOUT"\} &
\{"type":"res","code":1000,"msg":"OK"\} \\
\midrule

\normalfont Thông tin phiên \newline \cmd{WHOAMI} &
\{"cmd":"WHOAMI"\} &
\{"type":"res","code":1000,"msg":"OK",\newline "data":\{"username":"son","role":"PLAYER","team\_id":5,"is\_captain":true,"state":"IDLE"\}\} \\
\midrule

\normalfont Đổi mật khẩu \newline \cmd{CHANGE\_PASSWORD} &
\{"cmd":"CHANGE\_PASSWORD",\newline "old\_password":"Abc123","new\_password":"Xyz789"\} &
\{"type":"res","code":1000,"msg":"OK"\} \\
\midrule

\normalfont Tạo đội \newline \cmd{TEAM\_CREATE} &
\{"cmd":"TEAM\_CREATE","name":"Rong Bay"\} &
\{"type":"res","code":1000,"msg":"OK",\newline "data":\{"team\_id":5\}\} \\
\midrule

\normalfont Xóa đội \newline \cmd{TEAM\_DELETE} &
\{"cmd":"TEAM\_DELETE","team\_id":5\} &
\{"type":"res","code":1000,"msg":"OK"\} \\
\midrule

\normalfont Danh sách đội \newline \cmd{TEAM\_LIST} &
\{"cmd":"TEAM\_LIST","offset":0,"limit":20\} &
\{"type":"res","code":1000,"msg":"OK",\newline "data":\{"teams":[\{"team\_id":5,"name":"Rong Bay","captain":"son","size":2,"state":"IDLE"\}],"has\_more":false,"next\_offset":null\}\} \\
\midrule

\normalfont Thành viên đội \newline \cmd{TEAM\_MEMBERS} &
\{"cmd":"TEAM\_MEMBERS","team\_id":5\} &
\{"type":"res","code":1000,"msg":"OK",\newline "data":\{"members":[\{"slot":0,"username":"son","is\_captain":true,"online":true\}]\}\} \\
\midrule

\normalfont Xin vào đội \newline \cmd{TEAM\_JOIN\_REQUEST} &
\{"cmd":"TEAM\_JOIN\_REQUEST","team\_id":5\} &
\{"type":"res","code":1000,"msg":"OK"\} \\
\midrule

\normalfont Danh sách xin vào \newline \cmd{TEAM\_JOIN\_LIST} &
\{"cmd":"TEAM\_JOIN\_LIST"\} &
\{"type":"res","code":1000,"msg":"OK",\newline "data":\{"requests":["huy","lan"]\}\} \\
\midrule

\normalfont Duyệt xin vào \newline \cmd{TEAM\_JOIN\_REPLY} &
\{"cmd":"TEAM\_JOIN\_REPLY",\newline "username":"huy","accept":true\} &
\{"type":"res","code":1000,"msg":"OK"\} \\
\midrule

\normalfont Mời vào đội \newline \cmd{TEAM\_INVITE} &
\{"cmd":"TEAM\_INVITE","username":"lan"\} &
\{"type":"res","code":1000,"msg":"OK"\} \\
\midrule

\normalfont Danh sách lời mời \newline \cmd{TEAM\_INVITE\_LIST} &
\{"cmd":"TEAM\_INVITE\_LIST"\} &
\{"type":"res","code":1000,"msg":"OK",\newline "data":\{"invites":[\{"team\_id":5,"team\_name":"Rong Bay","from":"son"\}]\}\} \\
\midrule

\normalfont Trả lời lời mời \newline \cmd{TEAM\_INVITE\_REPLY} &
\{"cmd":"TEAM\_INVITE\_REPLY","team\_id":5,"accept":true\} &
\{"type":"res","code":1000,"msg":"OK"\} \\
\midrule

\normalfont Rời đội \newline \cmd{TEAM\_LEAVE} &
\{"cmd":"TEAM\_LEAVE"\} &
\{"type":"res","code":1000,"msg":"OK"\} \\
\midrule

\normalfont Xóa thành viên \newline \cmd{TEAM\_KICK} &
\{"cmd":"TEAM\_KICK","username":"lan"\} &
\{"type":"res","code":1000,"msg":"OK"\} \\
\midrule

\normalfont Chuyển đội trưởng \newline \cmd{TEAM\_TRANSFER} &
\{"cmd":"TEAM\_TRANSFER","username":"huy"\} &
\{"type":"res","code":1000,"msg":"OK"\} \\
\midrule

\normalfont Gửi thách đấu \newline \cmd{CHALLENGE\_SEND} &
\{"cmd":"CHALLENGE\_SEND","team\_id":8\} &
\{"type":"res","code":1000,"msg":"OK"\} \\
\midrule

\normalfont Danh sách thách đấu \newline \cmd{CHALLENGE\_LIST} &
\{"cmd":"CHALLENGE\_LIST"\} &
\{"type":"res","code":1000,"msg":"OK",\newline "data":\{"incoming":[\{"team\_id":8,"team\_name":"Sao Bang"\}],"outgoing":[]\}\} \\
\midrule

\normalfont Trả lời thách đấu \newline \cmd{CHALLENGE\_REPLY} &
\{"cmd":"CHALLENGE\_REPLY","team\_id":8,"accept":true\} &
\{"type":"res","code":1000,"msg":"OK",\newline "data":\{"lobby\_id":3\}\} \\
\midrule

\normalfont Hủy thách đấu \newline \cmd{CHALLENGE\_CANCEL} &
\{"cmd":"CHALLENGE\_CANCEL","team\_id":8\} &
\{"type":"res","code":1000,"msg":"OK"\} \\
\midrule

\normalfont Cài đặt sảnh \newline \cmd{LOBBY\_SETTINGS} &
\{"cmd":"LOBBY\_SETTINGS",\newline "turn\_time":30,"chest\_time":30,"start\_coin":0\} &
\{"type":"res","code":1000,"msg":"OK"\} \\
\midrule

\normalfont Sẵn sàng \newline \cmd{LOBBY\_READY} &
\{"cmd":"LOBBY\_READY","ready":true\} &
\{"type":"res","code":1000,"msg":"OK"\} \\
\midrule

\normalfont Bắt đầu trận \newline \cmd{LOBBY\_START} &
\{"cmd":"LOBBY\_START"\} &
\{"type":"res","code":1000,"msg":"OK"\} \\
\midrule

\normalfont Rời sảnh \newline \cmd{LOBBY\_LEAVE} &
\{"cmd":"LOBBY\_LEAVE"\} &
\{"type":"res","code":1000,"msg":"OK"\} \\
\midrule

\normalfont Lấy trạng thái trận \newline \cmd{GAME\_STATE} &
\{"cmd":"GAME\_STATE"\} &
\{"type":"res","code":1000,"msg":"OK",\newline "data":\{"state":\{...\}\}\} \\
\midrule

\normalfont Mua vật phẩm \newline \cmd{BUY} &
\{"cmd":"BUY","item":"LASER","qty":1\} &
\{"type":"res","code":1000,"msg":"OK"\} \\
\midrule

\normalfont Sửa tàu \newline \cmd{REPAIR} &
\{"cmd":"REPAIR","hp":150\} &
\{"type":"res","code":1000,"msg":"OK"\} \\
\midrule

\normalfont Nạp đạn \newline \cmd{RELOAD} &
\{"cmd":"RELOAD","weapon":"LASER","index":0\} &
\{"type":"res","code":1000,"msg":"OK"\} \\
\midrule

\normalfont Mở rương \newline \cmd{CHEST\_OPEN} &
\{"cmd":"CHEST\_OPEN"\} &
\{"type":"res","code":1000,"msg":"OK"\} \\
\midrule

\normalfont Trả lời rương \newline \cmd{CHEST\_ANSWER} &
\{"cmd":"CHEST\_ANSWER","choice":"B"\} &
\{"type":"res","code":1000,"msg":"OK"\} \\
\midrule

\normalfont Tấn công \newline \cmd{ATTACK} &
\{"cmd":"ATTACK","target":"minh",\newline "gun":\{"type":"LASER","count":2\},"missiles":1\} &
\{"type":"res","code":1000,"msg":"OK"\} \\
\midrule

\normalfont Danh sách người dùng \newline \cmd{ADMIN\_USER\_LIST} &
\{"cmd":"ADMIN\_USER\_LIST","offset":0,"limit":20\} &
\{"type":"res","code":1000,"msg":"OK",\newline "data":\{"users":[\{"username":"son","role":"PLAYER","locked":false,"online":true\}],"has\_more":false,"next\_offset":null\}\} \\
\midrule

\normalfont Khóa/mở tài khoản \newline \cmd{ADMIN\_LOCK} &
\{"cmd":"ADMIN\_LOCK","username":"son","locked":true\} &
\{"type":"res","code":1000,"msg":"OK"\} \\
\midrule

\normalfont Đọc log \newline \cmd{ADMIN\_LOG} &
\{"cmd":"ADMIN\_LOG","offset":0,"limit":20\} &
\{"type":"res","code":1000,"msg":"OK",\newline "data":\{"lines":["2026-11-18 INFO AUTH LOGIN user=son"],"has\_more":false,"next\_offset":null\}\} \\
\midrule

\normalfont Danh sách câu hỏi \newline \cmd{ADMIN\_QUESTION\_LIST} &
\{"cmd":"ADMIN\_QUESTION\_LIST","offset":0,"limit":20\} &
\{"type":"res","code":1000,"msg":"OK",\newline "data":\{"questions":[\{"question\_id":7,"tier":"BRONZE","question":"2+2=?","choices":["3","4","5","6"],"answer":"B"\}],"has\_more":false,"next\_offset":null\}\} \\
\midrule

\normalfont Thêm câu hỏi \newline \cmd{ADMIN\_QUESTION\_ADD} &
\{"cmd":"ADMIN\_QUESTION\_ADD",\newline "tier":"BRONZE","question":"2+2=?",\newline "choices":["3","4","5","6"],"answer":"B"\} &
\{"type":"res","code":1000,"msg":"OK",\newline "data":\{"question\_id":7\}\} \\
\midrule

\normalfont Xóa câu hỏi \newline \cmd{ADMIN\_QUESTION\_DELETE} &
\{"cmd":"ADMIN\_QUESTION\_DELETE","question\_id":7\} &
\{"type":"res","code":1000,"msg":"OK"\} \\
\bottomrule
\caption{Request và response message của toàn bộ lệnh}\label{tab:commands-detail}
\end{longtable}

Trong bảng chi tiết lệnh, response thành công dùng mã cụ thể \cmd{1000}. Các ký hiệu \cmd{1xxx}, \cmd{2xxx}, \ldots{} ở bảng nhóm mã chỉ biểu diễn từng nhóm kết quả; response lỗi giữ nguyên envelope, dùng một mã cụ thể thuộc nhóm tương ứng và thường không có \cmd{data}. \cmd{BUY.item} nhận một trong các giá trị \cmd{AMMO\_BOX}, \cmd{LASER}, \cmd{POWER\_PACK}, \cmd{MISSILE}, \cmd{ARMOR\_BASIC}, \cmd{ARMOR\_ADVANCED}; riêng hai loại giáp yêu cầu \cmd{qty=1}. \cmd{tier} nhận \cmd{"BRONZE"}, \cmd{"SILVER"} hoặc \cmd{"GOLD"}. Với \cmd{ATTACK}, phải có ít nhất \cmd{gun} hoặc \cmd{missiles>0}; \cmd{AUTO} luôn có \cmd{count=1}.

\subsection{Thứ tự response và event}

Mỗi kết nối chỉ được có \textbf{một request đang chờ response}. Client khóa thao tác hoặc xếp request mới vào hàng đợi cục bộ cho tới khi nhận response trước; server cũng xử lý request trên cùng một kết nối theo đúng thứ tự nhận. Mỗi request nhận đúng một response; response được đưa vào buffer gửi \textbf{trước} mọi event do chính request đó tạo ra. Cách này giúp client bỏ trạng thái ``đang chờ'' trước khi cập nhật giao diện từ event mà không cần trường \cmd{id}.

\begin{itemize}
  \item \cmd{ATTACK} thành công: response \cmd{ATTACK} $\rightarrow$ event \cmd{ATTACK} $\rightarrow$ \cmd{SHIP\_DESTROYED} nếu có $\rightarrow$ \cmd{MATCH\_END} hoặc \cmd{TURN\_START}.
  \item \cmd{TEAM\_JOIN\_REPLY} thành công: response cho đội trưởng $\rightarrow$ \cmd{JOIN\_RESULT} $\rightarrow$ \cmd{TEAM\_UPDATE}.
  \item \cmd{LOGIN} khi quay lại trận: response \cmd{LOGIN} $\rightarrow$ \cmd{MATCH\_RESUME} cho người đăng nhập $\rightarrow$ \cmd{PLAYER\_RECONNECTED} cho những người trong trận.
  \item Event do kết nối khác hoặc timer tạo ra có thể đến trong lúc client đang chờ response. Client phân biệt bằng \cmd{type}: \cmd{"res"} thuộc request duy nhất đang chờ, còn \cmd{"evt"} được chuyển cho bộ xử lý event.
\end{itemize}

\subsection{Chi tiết event}

Event là thông báo chủ động từ server nên không có request tương ứng. Mọi event đều có dạng \cmd{type="evt"}, \cmd{evt:string} và \cmd{data:object}. Trong các schema dưới đây, \cmd{member} và \cmd{settings} dùng kiểu đã định nghĩa ở mục trước; \cmd{team\_summary:object} gồm \cmd{team\_id:int}, \cmd{name:string} và \cmd{members:array<member>}.

\cmd{state:object} là snapshot trận được tạo riêng cho từng người nhận, gồm:
\begin{itemize}
  \item \cmd{match\_id:int}; \cmd{turn:object} gồm \cmd{player:string}, \cmd{deadline\_ms:int}, \cmd{paused:bool}.
  \item \cmd{chest:object|null}; nếu khác \cmd{null} thì gồm \cmd{tier:string} và \cmd{opened:bool}.
  \item \cmd{teams:array<object>}; mỗi đội gồm \cmd{team\_id:int}, \cmd{name:string}, \cmd{ships:array<ship>}.
  \item \cmd{ship} luôn có \cmd{username:string}, \cmd{slot:int}, \cmd{online:bool}, \cmd{alive:bool}, \cmd{hp:int}, \cmd{armor:array<object>} và \cmd{missiles:int}. Mỗi lớp giáp gồm \cmd{type:string}, \cmd{amor:int}.
  \item Với tàu cùng đội, \cmd{ship} có thêm \cmd{coin:int}, \cmd{auto:object} gồm \cmd{mag:int}, \cmd{boxes:int}; \cmd{lasers:array<object>} với mỗi phần tử có \cmd{mag:int}; và \cmd{power\_packs:int}. Với đối phương, các trường này bị ẩn và thay bằng \cmd{laser\_count:int}.
\end{itemize}

\begin{longtable}{@{}P{0.24\linewidth} J{0.68\linewidth}@{}}
\toprule
\textbf{Sự kiện} & \multicolumn{1}{P{0.68\linewidth}}{\textbf{Event message}} \\
\midrule\endfirsthead
\toprule
\textbf{Sự kiện} & \multicolumn{1}{P{0.68\linewidth}}{\textbf{Event message}} \\
\midrule\endhead

\normalfont Buộc đăng xuất \newline \cmd{FORCE\_LOGOUT} &
\{"type":"evt","evt":"FORCE\_LOGOUT",\newline "data":\{"reason":"LOCKED"\}\} \\
\midrule

\normalfont Có người xin vào \newline \cmd{JOIN\_REQUEST} &
\{"type":"evt","evt":"JOIN\_REQUEST",\newline "data":\{"username":"huy"\}\} \\
\midrule

\normalfont Kết quả xin vào \newline \cmd{JOIN\_RESULT} &
\{"type":"evt","evt":"JOIN\_RESULT",\newline "data":\{"team\_id":5,"accepted":true\}\} \\
\midrule

\normalfont Nhận lời mời \newline \cmd{INVITE} &
\{"type":"evt","evt":"INVITE",\newline "data":\{"team\_id":5,"team\_name":"Rong Bay","from":"son"\}\} \\
\midrule

\normalfont Kết quả lời mời \newline \cmd{INVITE\_RESULT} &
\{"type":"evt","evt":"INVITE\_RESULT",\newline "data":\{"username":"huy","accepted":true\}\} \\
\midrule

\normalfont Cập nhật đội \newline \cmd{TEAM\_UPDATE} &
\{"type":"evt","evt":"TEAM\_UPDATE",\newline "data":\{"team\_id":5,"members":[\{"slot":0,"username":"son","is\_captain":true,"online":true\}]\}\} \\
\midrule

\normalfont Bị xóa khỏi đội \newline \cmd{KICKED} &
\{"type":"evt","evt":"KICKED",\newline "data":\{"team\_id":5,"by":"son"\}\} \\
\midrule

\normalfont Đội bị xóa \newline \cmd{TEAM\_DELETED} &
\{"type":"evt","evt":"TEAM\_DELETED",\newline "data":\{"team\_id":5,"by":"son"\}\} \\
\midrule

\normalfont Nhận thách đấu \newline \cmd{CHALLENGE} &
\{"type":"evt","evt":"CHALLENGE",\newline "data":\{"team\_id":5,"team\_name":"Rong Bay"\}\} \\
\midrule

\normalfont Kết quả thách đấu \newline \cmd{CHALLENGE\_RESULT} &
\{"type":"evt","evt":"CHALLENGE\_RESULT",\newline "data":\{"team\_id":8,"accepted":true\}\} \\
\midrule

\normalfont Thách đấu bị hủy \newline \cmd{CHALLENGE\_CANCELLED} &
\{"type":"evt","evt":"CHALLENGE\_CANCELLED",\newline "data":\{"team\_id":5\}\} \\
\midrule

\normalfont Cập nhật sảnh \newline \cmd{LOBBY\_UPDATE} &
\{"type":"evt","evt":"LOBBY\_UPDATE",\newline "data":\{"lobby\_id":3,"host":\{"team\_id":5,"name":"Rong Bay","members":[...]\},"guest":\{"team\_id":8,"name":"Sao Bang","members":[...]\},"settings":\{"turn\_time":30,"chest\_time":30,"start\_coin":0\},"guest\_ready":true\}\} \\
\midrule

\normalfont Đóng sảnh \newline \cmd{LOBBY\_CLOSED} &
\{"type":"evt","evt":"LOBBY\_CLOSED",\newline "data":\{"reason":"DISCONNECTED"\}\} \\
\midrule

\normalfont Bắt đầu trận \newline \cmd{MATCH\_START} &
\{"type":"evt","evt":"MATCH\_START",\newline "data":\{"match\_id":3,"settings":\{"turn\_time":30,"chest\_time":30,"start\_coin":0\},"state":\{...\}\}\} \\
\midrule

\normalfont Tiếp tục trận \newline \cmd{MATCH\_RESUME} &
\{"type":"evt","evt":"MATCH\_RESUME",\newline "data":\{"match\_id":3,"settings":\{"turn\_time":30,"chest\_time":30,"start\_coin":0\},"state":\{...\}\}\} \\
\midrule

\normalfont Bắt đầu lượt \newline \cmd{TURN\_START} &
\{"type":"evt","evt":"TURN\_START",\newline "data":\{"player":"son","deadline\_ms":1763430146000,"state":\{...\}\}\} \\
\midrule

\normalfont Thả rương \newline \cmd{CHEST\_DROP} &
\{"type":"evt","evt":"CHEST\_DROP",\newline "data":\{"player":"son","tier":"SILVER"\}\} \\
\midrule

\normalfont Câu hỏi rương \newline \cmd{CHEST\_QUESTION} &
\{"type":"evt","evt":"CHEST\_QUESTION",\newline "data":\{"tier":"SILVER","question":"2+2=?","choices":["3","4","5","6"],"deadline\_ms":1763430150000,"can\_answer":true\}\} \\
\midrule

\normalfont Kết quả rương \newline \cmd{CHEST\_RESULT} &
\{"type":"evt","evt":"CHEST\_RESULT",\newline "data":\{"correct":true,"answer":"B","coin":500,"state":\{...\}\}\} \\
\midrule

\normalfont Mua vật phẩm \newline \cmd{PURCHASE} &
\{"type":"evt","evt":"PURCHASE",\newline "data":\{"player":"son","item":"LASER","qty":1,"state":\{...\}\}\} \\
\midrule

\normalfont Sửa tàu \newline \cmd{REPAIR} &
\{"type":"evt","evt":"REPAIR",\newline "data":\{"player":"son","hp":150,"state":\{...\}\}\} \\
\midrule

\normalfont Nạp đạn \newline \cmd{RELOAD} &
\{"type":"evt","evt":"RELOAD",\newline "data":\{"player":"son","weapon":"LASER","index":0,"state":\{...\}\}\} \\
\midrule

\normalfont Tấn công \newline \cmd{ATTACK} &
\{"type":"evt","evt":"ATTACK",\newline "data":\{"attacker":"son","target":"minh","gun":\{"type":"LASER","count":2\},"missiles":1,"damage":1000,"armor\_lost":500,"hp\_lost":500,"destroyed":false,"state":\{...\}\}\} \\
\midrule

\normalfont Bỏ lượt \newline \cmd{TURN\_SKIPPED} &
\{"type":"evt","evt":"TURN\_SKIPPED",\newline "data":\{"player":"son","reason":"TIMEOUT","state":\{...\}\}\} \\
\midrule

\normalfont Mất kết nối \newline \cmd{PLAYER\_DISCONNECTED} &
\{"type":"evt","evt":"PLAYER\_DISCONNECTED",\newline "data":\{"player":"son","state":\{...\}\}\} \\
\midrule

\normalfont Kết nối lại \newline \cmd{PLAYER\_RECONNECTED} &
\{"type":"evt","evt":"PLAYER\_RECONNECTED",\newline "data":\{"player":"son","state":\{...\}\}\} \\
\midrule

\normalfont Tàu bị hạ \newline \cmd{SHIP\_DESTROYED} &
\{"type":"evt","evt":"SHIP\_DESTROYED",\newline "data":\{"player":"minh","reason":"ATTACK","state":\{...\}\}\} \\
\midrule

\normalfont Kết thúc trận \newline \cmd{MATCH\_END} &
\{"type":"evt","evt":"MATCH\_END",\newline "data":\{"winner\_team\_id":5,"loser\_team\_id":8,"reason":"ELIMINATED","turns":12,"stats":[\{"username":"son","team\_id":5,"survived":true,"damage\_dealt":1800\}],"state":\{...\}\}\} \\
\bottomrule
\caption{Event message server gửi cho client}\label{tab:events-detail}
\end{longtable}

Snapshot trận được làm riêng cho từng người nhận: đội mình thấy đầy đủ; tàu đối phương chỉ hiện HP, giáp, số laze và số tên lửa, \emph{không} có coin và đạn.

% =====================================================================
\section{Responsibility}

Khối lượng công việc được phân chia cân bằng, mỗi thành viên chịu trách nhiệm chính cho 50\% công việc và cùng phối hợp tích hợp, kiểm thử toàn bộ hệ thống.

\begin{table}[H]\centering\small
\begin{tabularx}{\linewidth}{@{}P{0.28\linewidth} L@{}}
\toprule
\textbf{Thành viên} & \textbf{Nhiệm vụ chính} \\
\midrule
Đặng Hoàng Sơn -- 20235994 & Thiết kế kiến trúc server và giao thức; xây dựng tầng kết nối TCP/epoll, quản lý phiên và xác thực; phát triển chức năng đội, thách đấu và sảnh chờ; viết unit test cho các module phụ trách. \\
\midrule
Ngô Đức Huy -- 20235946 & Xây dựng luật chơi và xử lý trận đấu, lượt chơi, rương và trang bị; phát triển client GTK và đồng bộ trạng thái; xây dựng lưu trữ SQLite, giao diện quản trị; viết unit test cho các module phụ trách. \\
\bottomrule
\end{tabularx}
\caption{Phân công trách nhiệm của các thành viên}
\end{table}

Hai thành viên cùng thực hiện kiểm thử tích hợp, sửa lỗi, hoàn thiện tài liệu và chuẩn bị bản chạy cuối cùng.

\end{document}
